\chapter{Algorithmique : variables et structures de contrôle}

\textit{L'algorithmique est la science de la résolution de problèmes par des suites d'instructions ordonnées. Ce chapitre présente les concepts fondamentaux pour écrire et comprendre des algorithmes : variables, types de données, opérations d'entrée-sortie, et les structures qui permettent de contrôler le déroulement du programme (conditions et boucles).}

\section{L'essentiel du cours}

\subsection{Notion d'algorithme}
Un \textbf{algorithme} est une suite finie et ordonnée d'instructions qui, lorsqu'elles sont exécutées, permettent de résoudre un problème donné. Il se compose de trois parties principales :
\begin{itemize}
    \item \textbf{Entrée} : données fournies à l'algorithme.
    \item \textbf{Traitement} : opérations effectuées sur les données.
    \item \textbf{Sortie} : résultats produits par l'algorithme.
\end{itemize}

\begin{definition}[Algorithme]
Un algorithme est une procédure pas à pas, définie de manière non ambiguë, qui transforme des données d'entrée en données de sortie en un nombre fini d'étapes.
\end{definition}

\subsection{Variables et types}
Une \textbf{variable} est un espace mémoire nommé qui peut contenir une valeur. Chaque variable a un \textbf{type} qui détermine la nature des données qu'elle peut stocker.

\begin{propriete}[Types de base]
\begin{itemize}
    \item \textbf{Entier} : nombres entiers (ex: 5, -12, 0)
    \item \textbf{Réel} : nombres à virgule (ex: 3.14, -2.5)
    \item \textbf{Caractère} : un seul caractère (ex: 'A', '7')
    \item \textbf{Chaîne} : suite de caractères (ex: "Bonjour")
    \item \textbf{Logique (Booléen)} : vrai ou faux (Vrai/Faux)
\end{itemize}
\end{propriete}

\subsection{Affectation, lecture, écriture}

\begin{aretenir}[Instructions de base]
\begin{itemize}
    \item \textbf{Affectation} : \texttt{Variable $\leftarrow$ Expression} — stocke le résultat de l'expression dans la variable.
    \item \textbf{Lecture} : \texttt{Lire(Variable)} — permet à l'utilisateur de saisir une valeur.
    \item \textbf{Écriture} : \texttt{Écrire(Expression)} — affiche le résultat à l'écran.
\end{itemize}
\end{aretenir}

\begin{exemple}[Exemple simple]
\begin{algo}
Algorithme CalculSomme
Variables
    a, b, somme : Entier
Début
    Écrire("Entrez deux nombres :")
    Lire(a)
    Lire(b)
    somme $\leftarrow$ a + b
    Écrire("La somme est : ", somme)
Fin
\end{algo}
\end{exemple}

\subsection{Structures conditionnelles}
Les structures conditionnelles permettent d'exécuter des instructions seulement si une condition est vraie.

\begin{definition}[Structure Si...Alors...Sinon]
\begin{algo}
Si (condition) Alors
    Instructions_si_vrai
Sinon
    Instructions_si_faux
FinSi
\end{algo}
\end{definition}

\begin{aretenir}[Opérateurs de comparaison]
\begin{itemize}
    \item \texttt{=} : égal à
    \item \texttt{<>} ou \texttt{$\neq$} : différent de
    \item \texttt{<}, \texttt{>}, \texttt{<=}, \texttt{>=} : comparaisons classiques
    \item \texttt{ET}, \texttt{OU}, \texttt{NON} : opérateurs logiques
\end{itemize}
\end{aretenir}

\subsection{Structures itératives (boucles)}
Les boucles permettent de répéter un bloc d'instructions plusieurs fois.

\begin{definition}[Boucle Pour]
Utilisée quand le nombre d'itérations est connu à l'avance.
\begin{algo}
Pour i $\leftarrow$ debut À fin [Pas pas] Faire
    Instructions
FinPour
\end{algo}
\end{definition}

\begin{definition}[Boucle Tant que]
Utilisée quand on répète tant qu'une condition est vraie (test en début).
\begin{algo}
Tant que (condition) Faire
    Instructions
FinTantQue
\end{algo}
\end{definition}

\begin{definition}[Boucle Répéter...Jusqu'à]
Utilisée quand on répète jusqu'à ce qu'une condition devienne vraie (test en fin, exécutée au moins une fois).
\begin{algo}
Répéter
    Instructions
Jusqu'à (condition)
\end{algo}
\end{definition}

\begin{aretenir}[Points clés à retenir]
\begin{itemize}
    \item Un algorithme commence toujours par \texttt{Algorithme Nom} et se termine par \texttt{Fin}.
    \item Les variables doivent être déclarées dans la section \texttt{Variables}.
    \item L'affectation se note avec une flèche $\leftarrow$.
    \item Les conditions dans les \texttt{Si} et les boucles doivent être entre parenthèses.
    \item La boucle \texttt{Répéter...Jusqu'à} s'exécute au moins une fois, contrairement à \texttt{Tant que}.
    \item Pour dérouler un algorithme, on trace un tableau de valeurs des variables à chaque étape.
\end{itemize}
\end{aretenir}

\section{Exercices}

\rubrique{Notions de base : variables et affectation}

\exo{}\diff{1}
Écrire un algorithme qui demande deux nombres à l'utilisateur, calcule leur produit et affiche le résultat.

\exo{}\diff{1}
Écrire un algorithme qui échange le contenu de deux variables A et B (sans utiliser de troisième variable).

\exo{}\diff{1}
Quel est l'affichage produit par l'algorithme suivant ?
\begin{algo}
Algorithme Test1
Variables
    x, y : Entier
Début
    x $\leftarrow$ 5
    y $\leftarrow$ 3
    x $\leftarrow$ x + y
    y $\leftarrow$ x - y
    x $\leftarrow$ x - y
    Écrire(x, y)
Fin
\end{algo}

\exo{}\diff{1}
Écrire un algorithme qui lit le nom et l'âge d'une personne, puis affiche un message du type : "Bonjour [nom], vous avez [âge] ans."

\rubrique{Structures conditionnelles}

\exo{}\diff{1}
Écrire un algorithme qui lit un nombre et affiche s'il est positif, négatif ou nul.

\exo{}\diff{1}
Écrire un algorithme qui lit trois nombres et affiche le plus grand des trois.

\exo{}\diff{2}
Écrire un algorithme qui détermine si une année saisie est bissextile. Une année est bissextile si elle est divisible par 4 mais pas par 100, ou si elle est divisible par 400.

\exo{}\diff{2}
Un magasin de Yaoundé applique une remise de 10\% si le montant des achats dépasse 50 000 FCFA. Écrire un algorithme qui lit le montant des achats et affiche le montant à payer après remise éventuelle.

\exo{}\diff{2}
Écrire un algorithme qui lit deux nombres et un opérateur (+, -, *, /) et effectue l'opération correspondante. Gérer le cas de la division par zéro.

\rubrique{Boucle Pour}

\exo{}\diff{1}
Écrire un algorithme qui affiche les nombres de 1 à 10.

\exo{}\diff{1}
Écrire un algorithme qui calcule la somme des N premiers entiers naturels (N saisi par l'utilisateur).

\exo{}\diff{2}
Écrire un algorithme qui lit 10 notes (sur 20) et affiche leur moyenne.

\exo{}\diff{2}
Écrire un algorithme qui affiche la table de multiplication d'un nombre saisi par l'utilisateur (de 1 à 10).

\exo{}\diff{2}
Écrire un algorithme qui calcule N! (factorielle de N), N étant saisi par l'utilisateur.

\rubrique{Boucle Tant que et Répéter...Jusqu'à}

\exo{}\diff{1}
Écrire un algorithme utilisant \texttt{Tant que} qui affiche les nombres pairs de 2 à 20.

\exo{}\diff{2}
Écrire un algorithme qui demande à l'utilisateur de saisir un mot de passe (supposé "CAMEROUN") et répète la demande tant que le mot de passe est incorrect.

\exo{}\diff{2}
Écrire un algorithme qui lit une suite de nombres positifs (la saisie s'arrête quand on entre un nombre négatif) et affiche leur somme.

\exo{}\diff{2}
Écrire un algorithme utilisant \texttt{Répéter...Jusqu'à} qui demande à l'utilisateur de deviner un nombre secret (fixé à 42 dans l'algorithme) et affiche "Trop grand" ou "Trop petit" à chaque tentative.

\exo{}\diff{3}
Écrire un algorithme qui calcule la moyenne d'une série de notes saisies par l'utilisateur. La saisie s'arrête lorsque l'utilisateur entre -1.

\rubrique{Déroulement et traçage d'algorithmes}

\exo{}\diff{2}
Dérouler l'algorithme suivant en complétant le tableau de valeurs pour les entrées a=7 et b=3.
\begin{algo}
Algorithme Division
Variables
    a, b, quotient, reste : Entier
Début
    quotient $\leftarrow$ 0
    reste $\leftarrow$ a
    Tant que (reste >= b) Faire
        reste $\leftarrow$ reste - b
        quotient $\leftarrow$ quotient + 1
    FinTantQue
    Écrire(quotient, reste)
Fin
\end{algo}

\exo{}\diff{2}
Dérouler l'algorithme suivant pour N=5.
\begin{algo}
Algorithme Suite
Variables
    N, i, u : Entier
Début
    Lire(N)
    u $\leftarrow$ 0
    Pour i $\leftarrow$ 1 À N Faire
        u $\leftarrow$ u + i
    FinPour
    Écrire(u)
Fin
\end{algo}

\exo{}\diff{3}
Dérouler l'algorithme suivant pour la valeur saisie 1234.
\begin{algo}
Algorithme Mystere
Variables
    n, chiffre, inverse : Entier
Début
    Lire(n)
    inverse $\leftarrow$ 0
    Tant que (n > 0) Faire
        chiffre $\leftarrow$ n MOD 10
        inverse $\leftarrow$ inverse * 10 + chiffre
        n $\leftarrow$ n DIV 10
    FinTantQue
    Écrire(inverse)
Fin
\end{algo}

\rubrique{Exercices type Baccalauréat}

\sujetbac[2020, Cameroun, C/D/E]
\exo{}\diff{3}
Écrire un algorithme qui lit les coefficients a, b, c d'une équation du second degré $ax^2 + bx + c = 0$ et affiche ses solutions (réelles ou complexes). On suppose que a $\neq$ 0.

\sujetbac[2021, Cameroun, C/D/E]
\exo{}\diff{3}
Écrire un algorithme qui détermine si un nombre entier N saisi est premier. Un nombre premier est un entier supérieur à 1 qui n'a que deux diviseurs : 1 et lui-même.

\sujetbac[2022, Cameroun, C/D/E]
\exo{}\diff{3}
Écrire un algorithme qui lit 10 notes d'élèves (sur 20) et affiche :
\begin{itemize}
    \item La moyenne de la classe.
    \item Le nombre d'élèves ayant une note supérieure ou égale à 10.
    \item La note maximale et la note minimale.
\end{itemize}

\sujetbac[2023, Cameroun, C/D/E]
\exo{}\diff{3}
Écrire un algorithme qui génère et affiche les N premiers termes de la suite de Fibonacci. La suite est définie par : $F_0 = 0$, $F_1 = 1$, et $F_n = F_{n-1} + F_{n-2}$ pour $n \geq 2$.

\exo{}\diff{3}
Écrire un algorithme qui lit un nombre entier positif et affiche sa représentation en binaire (sans utiliser de tableau).

\exo{}\diff{3}
Écrire un algorithme qui lit un tableau de 10 entiers et affiche le nombre d'occurrences de chaque valeur (entre 0 et 9). On suppose que les valeurs saisies sont comprises entre 0 et 9.

\section{Corrigés}

\corrige{de l'exercice 1}
\begin{algo}
Algorithme Produit
Variables
    a, b, produit : Entier
Début
    Écrire("Entrez deux nombres :")
    Lire(a)
    Lire(b)
    produit $\leftarrow$ a * b
    Écrire("Le produit est : ", produit)
Fin
\end{algo}

\corrige{de l'exercice 2}
Pour échanger deux variables sans troisième variable, on utilise les opérations arithmétiques :
\begin{algo}
Algorithme Echange
Variables
    A, B : Entier
Début
    Écrire("Entrez A et B :")
    Lire(A)
    Lire(B)
    A $\leftarrow$ A + B
    B $\leftleftarrow$ A - B
    A $\leftarrow$ A - B
    Écrire("Après échange : A = ", A, ", B = ", B)
Fin
\end{algo}
\textbf{Explication} : Après A $\leftarrow$ A+B, A contient la somme. B $\leftarrow$ A-B donne la valeur initiale de A. Enfin A $\leftarrow$ A-B donne la valeur initiale de B.

\corrige{de l'exercice 3}
Déroulons l'algorithme pas à pas :
\begin{itemize}
    \item Initialement : x=5, y=3
    \item x $\leftarrow$ x + y = 5 + 3 = 8
    \item y $\leftarrow$ x - y = 8 - 3 = 5
    \item x $\leftarrow$ x - y = 8 - 5 = 3
\end{itemize}
L'affichage est : \textbf{3 5} (les valeurs ont été échangées).

\corrige{de l'exercice 4}
\begin{algo}
Algorithme Presentation
Variables
    nom : Chaîne
    age : Entier
Début
    Écrire("Entrez votre nom :")
    Lire(nom)
    Écrire("Entrez votre âge :")
    Lire(age)
    Écrire("Bonjour ", nom, ", vous avez ", age, " ans.")
Fin
\end{algo}

\corrige{de l'exercice 5}
\begin{algo}
Algorithme SigneNombre
Variables
    n : Entier
Début
    Écrire("Entrez un nombre :")
    Lire(n)
    Si (n > 0) Alors
        Écrire("Le nombre est positif.")
    Sinon
        Si (n < 0) Alors
            Écrire("Le nombre est négatif.")
        Sinon
            Écrire("Le nombre est nul.")
        FinSi
    FinSi
Fin
\end{algo}

\corrige{de l'exercice 6}
\begin{algo}
Algorithme PlusGrand
Variables
    a, b, c, max : Entier
Début
    Écrire("Entrez trois nombres :")
    Lire(a)
    Lire(b)
    Lire(c)
    max $\leftarrow$ a
    Si (b > max) Alors
        max $\leftarrow$ b
    FinSi
    Si (c > max) Alors
        max $\leftarrow$ c
    FinSi
    Écrire("Le plus grand est : ", max)
Fin
\end{algo}

\corrige{de l'exercice 7}
\begin{algo}
Algorithme AnneeBissextile
Variables
    annee : Entier
Début
    Écrire("Entrez une année :")
    Lire(annee)
    Si ((annee MOD 4 = 0 ET annee MOD 100 <> 0) OU (annee MOD 400 = 0)) Alors
        Écrire(annee, " est bissextile.")
    Sinon
        Écrire(annee, " n'est pas bissextile.")
    FinSi
Fin
\end{algo}

\corrige{de l'exercice 8}
\begin{algo}
Algorithme Remise
Variables
    montant, aPayer : Réel
Début
    Écrire("Entrez le montant des achats (en FCFA) :")
    Lire(montant)
    Si (montant > 50000) Alors
        aPayer $\leftarrow$ montant * 0.9
    Sinon
        aPayer $\leftarrow$ montant
    FinSi
    Écrire("Montant à payer : ", aPayer, " FCFA")
Fin
\end{algo}

\corrige{de l'exercice 9}
\begin{algo}
Algorithme Calculatrice
Variables
    a, b, resultat : Réel
    operateur : Caractère
Début
    Écrire("Entrez deux nombres :")
    Lire(a)
    Lire(b)
    Écrire("Entrez l'opérateur (+, -, *, /) :")
    Lire(operateur)
    Selon operateur
        Cas '+' : resultat $\leftarrow$ a + b
        Cas '-' : resultat $\leftarrow$ a - b
        Cas '*' : resultat $\leftarrow$ a * b
        Cas '/' : 
            Si (b <> 0) Alors
                resultat $\leftarrow$ a / b
            Sinon
                Écrire("Erreur : division par zéro !")
                resultat $\leftarrow$ 0
            FinSi
        Sinon : Écrire("Opérateur invalide !")
    FinSelon
    Écrire("Résultat : ", resultat)
Fin
\end{algo}

\corrige{de l'exercice 10}
\begin{algo}
Algorithme Compter1a10
Variables
    i : Entier
Début
    Pour i $\leftarrow$ 1 À 10 Faire
        Écrire(i)
    FinPour
Fin
\end{algo}

\corrige{de l'exercice 11}
\begin{algo}
Algorithme SommeN
Variables
    N, i, somme : Entier
Début
    Écrire("Entrez un entier N :")
    Lire(N)
    somme $\leftarrow$ 0
    Pour i $\leftarrow$ 1 À N Faire
        somme $\leftarrow$ somme + i
    FinPour
    Écrire("La somme des ", N, " premiers entiers est : ", somme)
Fin
\end{algo}

\corrige{de l'exercice 12}
\begin{algo}
Algorithme Moyenne10Notes
Variables
    i : Entier
    note, somme, moyenne : Réel
Début
    somme $\leftarrow$ 0
    Pour i $\leftarrow$ 1 À 10 Faire
        Écrire("Entrez la note ", i, " :")
        Lire(note)
        somme $\leftarrow$ somme + note
    FinPour
    moyenne $\leftarrow$ somme / 10
    Écrire("La moyenne est : ", moyenne)
Fin
\end{algo}

\corrige{de l'exercice 13}
\begin{algo}
Algorithme TableMultiplication
Variables
    n, i : Entier
Début
    Écrire("Entrez un nombre :")
    Lire(n)
    Pour i $\leftarrow$ 1 À 10 Faire
        Écrire(n, " x ", i, " = ", n * i)
    FinPour
Fin
\end{algo}

\corrige{de l'exercice 14}
\begin{algo}
Algorithme Factorielle
Variables
    N, i, fact : Entier
Début
    Écrire("Entrez un entier N :")
    Lire(N)
    fact $\leftarrow$ 1
    Pour i $\leftarrow$ 2 À N Faire
        fact $\leftarrow$ fact * i
    FinPour
    Écrire(N, "! = ", fact)
Fin
\end{algo}

\corrige{de l'exercice 15}
\begin{algo}
Algorithme NombresPairs
Variables
    i : Entier
Début
    i $\leftarrow$ 2
    Tant que (i <= 20) Faire
        Écrire(i)
        i $\leftarrow$ i + 2
    FinTantQue
Fin
\end{algo}

\corrige{de l'exercice 16}
\begin{algo}
Algorithme MotDePasse
Variables
    mdp : Chaîne
Début
    Écrire("Entrez le mot de passe :")
    Lire(mdp)
    Tant que (mdp <> "CAMEROUN") Faire
        Écrire("Mot de passe incorrect. Réessayez :")
        Lire(mdp)
    FinTantQue
    Écrire("Accès autorisé.")
Fin
\end{algo}

\corrige{de l'exercice 17}
\begin{algo}
Algorithme SommePositifs
Variables
    n, somme : Entier
Début
    somme $\leftarrow$ 0
    Écrire("Entrez un nombre (négatif pour arrêter) :")
    Lire(n)
    Tant que (n >= 0) Faire
        somme $\leftarrow$ somme + n
        Écrire("Entrez un nombre (négatif pour arrêter) :")
        Lire(n)
    FinTantQue
    Écrire("La somme des nombres positifs est : ", somme)
Fin
\end{algo}

\corrige{de l'exercice 18}
\begin{algo}
Algorithme Devinette
Variables
    secret, essai : Entier
Début
    secret $\leftarrow$ 42
    Répéter
        Écrire("Devinez le nombre secret :")
        Lire(essai)
        Si (essai > secret) Alors
            Écrire("Trop grand !")
        Sinon
            Si (essai < secret) Alors
                Écrire("Trop petit !")
            FinSi
        FinSi
    Jusqu'à (essai = secret)
    Écrire("Bravo ! Vous avez trouvé le nombre secret.")
Fin
\end{algo}

\corrige{de l'exercice 19}
\begin{algo}
Algorithme MoyenneNotes
Variables
    note, somme, moyenne : Réel
    nbNotes : Entier
Début
    somme $\leftarrow$ 0
    nbNotes $\leftarrow$ 0
    Écrire("Entrez une note (-1 pour terminer) :")
    Lire(note)
    Tant que (note <> -1) Faire
        somme $\leftarrow$ somme + note
        nbNotes $\leftarrow$ nbNotes + 1
        Écrire("Entrez une note (-1 pour terminer) :")
        Lire(note)
    FinTantQue
    Si (nbNotes > 0) Alors
        moyenne $\leftarrow$ somme / nbNotes
        Écrire("La moyenne est : ", moyenne)
    Sinon
        Écrire("Aucune note saisie.")
    FinSi
Fin
\end{algo}

\corrige{de l'exercice 20}
Tableau de déroulement pour a=7, b=3 :

\begin{center}
\begin{tabular}{|c|c|c|c|c|}
\hline
\textbf{Étape} & \textbf{quotient} & \textbf{reste} & \textbf{Condition (reste >= b)} & \textbf{Action} \\
\hline
Initial & 0 & 7 & — & — \\
\hline
1 & 0 & 7 & 7 >= 3 : Vrai & reste $\leftarrow$ 7-3=4, quotient $\leftarrow$ 0+1=1 \\
\hline
2 & 1 & 4 & 4 >= 3 : Vrai & reste $\leftarrow$ 4-3=1, quotient $\leftarrow$ 1+1=2 \\
\hline
3 & 2 & 1 & 1 >= 3 : Faux & Sortie de la boucle \\
\hline
\end{tabular}
\end{center}
Affichage : \textbf{2 1} (quotient = 2, reste = 1).

\corrige{de l'exercice 21}
Tableau de déroulement pour N=5 :

\begin{center}
\begin{tabular}{|c|c|c|c|}
\hline
\textbf{i} & \textbf{u (avant)} & \textbf{u (après)} & \textbf{Commentaire} \\
\hline
Initial & 0 & — & u $\leftarrow$ 0 \\
\hline
1 & 0 & 1 & u $\leftarrow$ 0 + 1 \\
\hline
2 & 1 & 3 & u $\leftarrow$ 1 + 2 \\
\hline
3 & 3 & 6 & u $\leftarrow$ 3 + 3 \\
\hline
4 & 6 & 10 & u $\leftarrow$ 6 + 4 \\
\hline
5 & 10 & 15 & u $\leftarrow$ 10 + 5 \\
\hline
\end{tabular}
\end{center}
Affichage : \textbf{15} (c'est la somme des 5 premiers entiers : 1+2+3+4+5=15).

\corrige{de l'exercice 22}
Tableau de déroulement pour n=1234 :

\begin{center}
\begin{tabular}{|c|c|c|c|c|}
\hline
\textbf{Étape} & \textbf{n} & \textbf{chiffre} & \textbf{inverse} & \textbf{Condition (n > 0)} \\
\hline
Initial & 1234 & — & 0 & — \\
\hline
1 & 123 & 4 & 0*10+4=4 & 1234 > 0 : Vrai \\
\hline
2 & 12 & 3 & 4*10+3=43 & 123 > 0 : Vrai \\
\hline
3 & 1 & 2 & 43*10+2=432 & 12 > 0 : Vrai \\
\hline
4 & 0 & 1 & 432*10+1=4321 & 1 > 0 : Vrai \\
\hline
5 & 0 & — & — & 0 > 0 : Faux \\
\hline
\end{tabular}
\end{center}
Affichage : \textbf{4321} (l'inverse du nombre saisi).

\corrige{de l'exercice 23}
\begin{algo}
Algorithme EquationSecondDegre
Variables
    a, b, c, discriminant : Réel
    x1, x2, partieReelle, partieImaginaire : Réel
Début
    Écrire("Entrez les coefficients a, b, c :")
    Lire(a)
    Lire(b)
    Lire(c)
    discriminant $\leftarrow$ b*b - 4*a*c
    Si (discriminant > 0) Alors
        x1 $\leftarrow$ (-b + sqrt(discriminant)) / (2*a)
        x2 $\leftarrow$ (-b - sqrt(discriminant)) / (2*a)
        Écrire("Deux solutions réelles : ", x1, " et ", x2)
    Sinon
        Si (discriminant = 0) Alors
            x1 $\leftarrow$ -b / (2*a)
            Écrire("Une solution double : ", x1)
        Sinon
            partieReelle $\leftarrow$ -b / (2*a)
            partieImaginaire $\leftarrow$ sqrt(-discriminant) / (2*a)
            Écrire("Deux solutions complexes : ", partieReelle, " + i*", partieImaginaire, " et ", partieReelle, " - i*", partieImaginaire)
        FinSi
    FinSi
Fin
\end{algo}

\corrige{de l'exercice 24}
\begin{algo}
Algorithme NombrePremier
Variables
    N, i : Entier
    estPremier : Booléen
Début
    Écrire("Entrez un entier N :")
    Lire(N)
    Si (N <= 1) Alors
        estPremier $\leftarrow$ Faux
    Sinon
        estPremier $\leftarrow$ Vrai
        Pour i $\leftarrow$ 2 À N-1 Faire
            Si (N MOD i = 0) Alors
                estPremier $\leftarrow$ Faux
            FinSi
        FinPour
    FinSi
    Si (estPremier = Vrai) Alors
        Écrire(N, " est un nombre premier.")
    Sinon
        Écrire(N, " n'est pas un nombre premier.")
    FinSi
Fin
\end{algo}

\corrige{de l'exercice 25}
\begin{algo}
Algorithme StatistiquesNotes
Variables
    i : Entier
    note, somme, moyenne, max, min : Réel
    nbReussite : Entier
Début
    somme $\leftarrow$ 0
    nbReussite $\leftarrow$ 0
    max $\leftarrow$ 0
    min $\leftarrow$ 20
    Pour i $\leftarrow$ 1 À 10 Faire
        Écrire("Entrez la note ", i, " (sur 20) :")
        Lire(note)
        somme $\leftarrow$ somme + note
        Si (note >= 10) Alors
            nbReussite $\leftarrow$ nbReussite + 1
        FinSi
        Si (note > max) Alors
            max $\leftarrow$ note
        FinSi
        Si (note < min) Alors
            min $\leftarrow$ note
        FinSi
    FinPour
    moyenne $\leftarrow$ somme / 10
    Écrire("Moyenne de la classe : ", moyenne)
    Écrire("Nombre d'élèves ayant la moyenne : ", nbReussite)
    Écrire("Note maximale : ", max)
    Écrire("Note minimale : ", min)
Fin
\end{algo}

\corrige{de l'exercice 26}
\begin{algo}
Algorithme Fibonacci
Variables
    N, i, f0, f1, fn : Entier
Début
    Écrire("Entrez le nombre de termes N :")
    Lire(N)
    f0 $\leftarrow$ 0
    f1 $\leftarrow$ 1
    Si (N >= 1) Alors
        Écrire(f0)
    FinSi
    Si (N >= 2) Alors
        Écrire(f1)
    FinSi
    Pour i $\leftarrow$ 3 À N Faire
        fn $\leftarrow$ f0 + f1
        Écrire(fn)
        f0 $\leftarrow$ f1
        f1 $\leftarrow$ fn
    FinPour
Fin
\end{algo}

\corrige{de l'exercice 27}
\begin{algo}
Algorithme ConversionBinaire
Variables
    n, reste : Entier
    binaire : Chaîne
Début
    Écrire("Entrez un nombre entier positif :")
    Lire(n)
    binaire $\leftarrow$ ""
    Si (n = 0) Alors
        binaire $\leftarrow$ "0"
    Sinon
        Tant que (n > 0) Faire
            reste $\leftarrow$ n MOD 2
            Si (reste = 0) Alors
                binaire $\leftarrow$ "0" + binaire
            Sinon
                binaire $\leftarrow$ "1" + binaire
            FinSi
            n $\leftarrow$ n DIV 2
        FinTantQue
    FinSi
    Écrire("Représentation binaire : ", binaire)
Fin
\end{algo}

\corrige{de l'exercice 28}
\begin{algo}
Algorithme Occurrences
Variables
    i, valeur : Entier
    occurrences : Tableau[0..9] d'Entier
Début
    Pour i $\leftarrow$ 0 À 9 Faire
        occurrences[i] $\leftarrow$ 0
    FinPour
    Pour i $\leftarrow$ 1 À 10 Faire
        Écrire("Entrez une valeur entre 0 et 9 :")
        Lire(valeur)
        occurrences[valeur] $\leftarrow$ occurrences[valeur] + 1
    FinPour
    Pour i $\leftarrow$ 0 À 9 Faire
        Écrire("La valeur ", i, " apparaît ", occurrences[i], " fois.")
    FinPour
Fin
\end{algo}